\section{Методы синтеза речи и правила транскрибирования}

\subsection{Классификация методов синтеза речи (задание 3.1)}

Все существующие методы синтеза речевого сигнала делятся на два основных
класса~\cite{lobanov}: \emph{параметрические} и \emph{компиляционные}.

\emph{Параметрические} методы основаны на знаниях о том, как образуется
речь у человека и какими параметрами описывается речевой сигнал. На этих
знаниях строится модель акустических процессов речеобразования, и сигнал
генерируется подобно тому, как это происходит в речевом тракте. Модель
может опираться либо на артикуляторные процессы, либо на математическое
описание собственно акустических процессов (формантная модель, модель
источник--фильтр).

\emph{Компиляционные} методы, наоборот, игнорируют физику речеобразования:
для синтеза используются отрезки естественной речевой волны, из которых
затем собирается (компилируется) сигнал. В качестве единиц компиляции
применяются:
\begin{itemize}
	\item микроволны --- отрезки, вырезанные из записей отдельных фонем или
	аллофонов;
	\item аллофоны --- реализации фонемы в конкретном фонетическом
	окружении;
	\item дифоны --- всевозможные пары соседних звуков, что позволяет
	сохранить естественные переходы между ними.
\end{itemize}

Компиляционные методы обладают важными достоинствами: синтезированная
речь звучит близко к естественной, а при наличии достаточной базы можно
<<научить>> компьютер говорить голосом конкретного человека, что
используется в системах клонирования голоса~\cite{lobanov,cloning}. К
недостаткам относятся трудоёмкость разработки и необходимость хранить
большие объёмы памяти для баз голосов.

\subsection{Правила транскрибирования}

Основная идея компиляционных методов состоит в использовании сегментов
естественной речевой волны, поэтому перед синтезом орфографический текст
нужно преобразовать в последовательность фонем. Для этого служат правила
транскрибирования; в каждом языке они свои. Рассмотрим основные правила
русской речи~\cite{pecheneva,zinder}.

В русском языке десять гласных букв (У, О, А, Э, Ы, И, Ю, Ё, Я, Е), но
только шесть гласных фонем ([У], [О], [А], [Э], [Ы], [И]). Буквы Ю, Ё, Я,
Е являются <<лишними>> и преобразуются одним из двух способов.

\begin{enumerate}
	\item Буква Ю, Ё, Я или Е превращается в пару фонем [Й] + [У], [О],
	[А] или [Э] соответственно, если она стоит в начале слова, после
	мягкого или твёрдого знака, а также после гласной. Например:
	вьюга \(\to\) [в й у г а], объём \(\to\) [а б й о м],
	яма \(\to\) [й а м а], поёт \(\to\) [п а й о т].
	\item В остальных случаях (после согласной) эти буквы дают фонемы
	[У], [О], [А], [Э] и смягчают предшествующую согласную. Например:
	тюк \(\to\) [т' у к], лёд \(\to\) [л' о т], пять \(\to\) [п' а т'],
	семь \(\to\) [с' э м'].
\end{enumerate}

Кроме того, действуют правила редукции безударных гласных: безударная
Я может переходить в [И] (тянуть \(\to\) [т' и н у т']), безударная О ---
в [А] (молоко \(\to\) [м а л а к о]), безударная Е --- в [И]
(дерево \(\to\) [д' э р и в а]). Буква И после Ж и Ш переходит в [Ы]:
жить \(\to\) [ж ы т'], шить \(\to\) [ш ы т'].

Для согласных важнейшими являются следующие правила:
\begin{enumerate}
	\item [Щ] образуется из сочетаний СЧ, ЗЧ, ЖЧ (счастье, извозчик,
	перебежчик);
	\item [Ц] --- из Ц и сочетания ТЬС (купаться);
	\item оглушение: звонкая согласная перед глухой или в конце слова
	переходит в парную глухую (лёд \(\to\) [л' о т], буг \(\to\) [б у к]);
	\item озвончение: глухая перед звонкой переходит в звонкую
	(косьба \(\to\) [к а з' б а]);
	\item уподобление и смягчение парных согласных (сшить \(\to\)
	[ш ы т'], хвостик \(\to\) [х в о с' т' и к]).
\end{enumerate}
Есть и исключения: <<бог>> \(\to\) [б о х], <<что>> \(\to\) [ш т о]. Все
перечисленные правила хорошо формализуются и реализуются алгоритмически.

\subsection{Обобщённая структура системы синтеза речи}

Система автоматического синтеза речи включает два основных блока:
\begin{enumerate}
	\item анализ текста и его преобразование в последовательность фонем
	(нормализация текста, транскрибирование, расстановка ударений);
	\item собственно синтез --- формирование речевого сигнала по фонемной
	последовательности.
\end{enumerate}
Для компиляционного синтеза второй блок распадается на выбор единиц из
базы, их конкатенацию и просодическую модификацию. Подробная схема
синтезатора, смоделированного в данной работе, приведена в
разделе~\ref{sec:structure}.

\subsection{Используемое программное обеспечение}

Запись сигналов выполнялась в программе Audacity 4.0.1 с параметрами
16~кГц, моно, 16~бит, которая позволяет записывать сигнал с микрофона,
просматривать осциллограмму, выделять, вырезать и вставлять фрагменты,
изменять громкость и сохранять результат в файл \texttt{*.wav}.

Анализ и обработка выполнены на языке Python (NumPy, SciPy, soundfile,
Matplotlib). Реализованы:
\begin{itemize}
	\item сегментация по кратковременной энергии (кадры 10~мс, порог 8\%
	от максимума, объединение близких сегментов);
	\item оценивание частоты основного тона $F_0$ автокорреляционным
	методом с выбором первого сильного максимума АКФ, октавной
	коррекцией и разрывом контура на паузах;
	\item вырезание фонем и слогов, приведение длительностей и уровней,
	склейка с перекрёстным затуханием;
	\item сдвиг тона с сохранением формант (кепстральная оболочка спектра
	и фазовый вокодер), изменение длительности без изменения тона.
\end{itemize}
